% chapters/03/sub-cpmk.tex
\begin{subcpmkbox}
Bab ini merupakan materi persiapan menuju \SUBCPMK{1.1}, yaitu ``menjelaskan
struktur \istilah{game loop} (input, \istilah{update}, \istilah{render}) dan
peran manajemen waktu dalam permainan''. Sub-CPMK tersebut baru dituntaskan
pada Bab~IV. Bab ini menyiapkan seluruh perkakas dan kebiasaan kerja yang
diperlukan untuk mencapainya.

\vspace{0.4em}
\begin{tabularx}{\textwidth}{@{}L{3.0cm} Y@{}}
\toprule
\headrow \textbf{Kedudukan} & \textbf{Keterangan} \\
\midrule
Materi RPS & Materi 1: Pengantar Pemrograman Game dan Persiapan Python/Pygame \\
\addlinespace
Sub-CPMK didekati & \SUBCPMK{1.1} (ranah C2), dituntaskan pada Bab~IV \\
\addlinespace
Keluaran bab & Lingkungan kerja yang siap dipakai, program Pygame pertama yang
berjalan, dan repositori tugas dengan penyimpanan awal \\
\addlinespace
Asesmen & Praktikum 01 (unjuk kerja) dengan rubrik penilaian artefak \\
\bottomrule
\end{tabularx}
\end{subcpmkbox}
